% Содержимое подраздела 1.3 - Генеративные модели на основе вариационных автоэнкодеров


\section{Генеративные модели на основе вариационных автоэнкодеров}
\label{sec:vae}

Несмотря на успехи RNN, они обладают существенным недостатком: процесс генерации сложно контролировать. Отсутствует явный механизм, позволяющий целенаправленно создавать молекулы с заданными свойствами. Для решения этой проблемы были предложены модели на основе \textbf{вариационных автоэнкодеров (Variational Autoencoder, VAE)}. VAE состоит из двух частей: энкодера, который сжимает дискретное представление молекулы (например, SMILES) в непрерывный вектор в скрытом (латентном) пространстве, и декодера, который восстанавливает из этого вектора исходную молекулу. Главное преимущество VAE заключается в том, что его латентное пространство является непрерывным и структурированным, что открывает возможность для градиентной оптимизации свойств и интерполяции между известными структурами для создания новых~\cite{wigh2022review}.

Однако стандартные VAE сталкиваются с рядом проблем, главной из которых является \textbf{коллапс апостериорного распределения (posterior collapse)}. Этот эффект возникает, когда декодер учится игнорировать информацию из латентного пространства и генерирует выходные данные, опираясь только на входные условия, что приводит к потере разнообразия. Для решения этой проблемы в работе, представляющей архитектуру \textbf{PCF-VAE}~\cite{bhadwal2025pcf}, была предложена модифицированная функция потерь. В ней вводится динамический весовой коэффициент $\alpha$ для члена KL-дивергенции, который постепенно увеличивается в процессе обучения (метод <<отжига>>). Это позволяет модели на ранних этапах сфокусироваться на точности реконструкции, а затем плавно <<включать>> регуляризацию латентного пространства. В сочетании с использованием более робастной нотации GenSMILES, данный подход позволил достичь на бенчмарке MOSES высокой валидности (98.01\%), уникальности (100\%) и новизны (93.77\%), демонстрируя эффективное решение проблемы коллапса.

Другой важной задачей является создание моделей для условной генерации, то есть создания молекул с заранее определенными свойствами. В работе, описывающей \textbf{условный $\beta$-VAE}~\cite{richards2022conditional}, предложен элегантный метод <<явного обусловливания латентного пространства>> (Explicit Latent Conditioning, ELC). Вместо того чтобы центрировать априорное распределение в латентном пространстве вокруг нуля, модель центрирует его вокруг стандартизированного значения целевого свойства $\hat{c}$, то есть использует распределение $N(\hat{c}, I)$. Это заставляет модель организовывать латентное пространство в соответствии с целевым свойством, группируя молекулы со схожими значениями вместе. Данный подход, реализованный в гибридной архитектуре (Transformer-энкодер и GRU-декодер), показал рекордные результаты в задаче неограниченной оптимизации, сгенерировав молекулу с предсказанным значением penalized LogP (pLogP) равным \textbf{104.29}, что на порядок превосходило предыдущие достижения.

Несмотря на успехи строковых моделей, проблема генерации невалидных SMILES остается актуальной. Это привело к развитию графовых VAE, которые работают не с линейной строкой, а с более естественным графовым представлением молекулы. Одним из наиболее успешных примеров является \textbf{Junction Tree VAE (JT-VAE)}~\cite{kondratyev2023generative}. Эта модель оперирует на уровне химически осмысленных строительных блоков (колец и связей), что по построению гарантирует \textbf{100\% химическую валидность} генерируемых структур. Такой подход позволяет избежать стадии валидации и сосредоточиться на оптимизации свойств в хорошо структурированном латентном пространстве.
